\chapter{Dimension and volume}\label{chap:dim}

\section{Linear dimension}

Let $\spc{A}$ be an Alexandrov space.
We define its \index{linear dimension}\emph{linear dimension} \index{LinDim@$\LinDim$}$\LinDim \spc{A}$ as the least upper bound of integers $m$ such that
the Euclidean space $\EE^m$ is isometric to a subspace of the tangent space $\T_p\spc{A}$ at \textit{some} point $p\in \spc{A}$.

If not stated otherwise, dimension means linear dimension.
In \ref{sec:packing-numbers} and \ref{sec:all-dim}, we will show that the linear dimension of Alexandrov spaces coincides with all reasonable notions of dimension;
thereafter, we will use the notation \index{dim@$\dim$}$\dim\spc{A}$ for $\LinDim\spc{A}$.

\begin{thm}{(\textit{n}+1)-comparison}\label{thm:n+1}
Let $\spc{A}$ be an $\Alex0$ space.
Then for any finite set of points $p,x_1,\dots,x_n\in \spc{A}$, there exists a model configuration
$\tilde p,\tilde x_1,\dots,\tilde x_n\in \EE^m$ for some $m$ such that
\[|\tilde p-\tilde x_i|_{\EE^m}=| p- x_i|_{\spc{A}}
\quad\text{and}\quad
|\tilde x_i-\tilde x_j|_{\EE^m}\ge |x_i- x_j|_{\spc{A}}\]
for any $i$ and $j$.
Moreover, we can assume that $m\le \LinDim\spc{A}$.

If $\spc{A}$ is an $\Alex{\pm1}$ space, then the same holds if we replace $\EE^m$ by $\SSS^m$ or $\HH^m$, respectively; in the $\Alex{1}$ case one has to assume in addition that $\diam \spc{A}\le \pi$.
\end{thm}

\parit{Proof.}
By \ref{cor:euclid-subcone}, we can choose a point $p'$ arbitrarily close to $p$ so that
$\Lin_{p'}\ni \dir{p'}{x_i}$ for any $i$.
Let us identify $\EE^m$ with a subspace of $\Lin_{p'}$ spanned by $\dir{p'}{x_1},\dots,\dir{p'}{x_n}$.
Note that $m\le \LinDim\spc{A}$.

Set $\tilde p'=0\in \EE^m$ and $\tilde x_i=\dist{p'}{x_i}{}\cdot\dir{p'}{x_i}\in \EE^m$ for every $i$.
Note that
\[|\tilde p'-\tilde x_i|_{\EE^m}=| p'- x_i|_{\spc{A}}\]
for every $i$.
Applying the comparison $\mangle\hinge {p'}{x_i}{x_j}\ge \angk {p'}{x_i}{x_j}$, we get
\[|\tilde x_i-\tilde x_j|_{\EE^m}\ge |x_i- x_j|_{\spc{A}}\]
for all $i$ and $j$.
Passing to a limit configuration as $p'\to p$, we get the result.

If $\spc{A}$ is an $\Alex{-1}$ space, then the same argument gives
$\tilde p,\tilde x_1,\z\dots,\tilde x_n\in \EE^m$ such that
\[|\tilde p-\tilde x_i|_{\EE^m}=| p- x_i|_{\spc{A}}
\quad\text{and}\quad
\side\hinge{\tilde p}{\tilde x_i}{\tilde x_j}_{\HH^2}\ge |x_i- x_j|_{\spc{A}}.\]
It remains to observe that $\EE^m$ with the metric defined by
$\dist{\tilde v}{\tilde w}{}\z\df \side\hinge{\tilde p}{\tilde v}{\tilde w}_{\HH^2}$
is isometric to $\HH^m$.

The $\Alex{1}$ case is analogous, but one has to pull $x_i$ toward $p$
to ensure that $\dist{p}{x_i}{}<\pi$ for each $i$, and then pass to an additional limit.
\qeds

\begin{thm}[!]{Exercise}\label{ex:dim-max}
Let $\spc{A}$ be an $\Alex0$ space.
Assume that $\EE^m$ is isometric to a subspace of the tangent space $\T_p\spc{A}$ at \textit{some} point $p\in \spc{A}$ and $\LinDim \spc{A}\z= m<\infty$.
Show that $\T_p\spc{A}\iso \EE^m$.
\end{thm}


\begin{thm}[!]{Exercise}\label{ex:dim=1}
Show that a 1-dimensional Alexandrov space is homeomorphic to a 1-dimensional manifold, possibly with nonempty boundary.
\end{thm}


\begin{thm}{Exercise}\label{ex:resporka}
Let $\spc{A}$ be an $\Alex0$ space.
Show that $\LinDim \spc{A}\z\ge m$ if and only if there are $m+2$ points $p$, $a_0,\z\dots, a_{m}\in \spc{A}$
such that $\angk p{a_i}{a_j}>\tfrac\pi2$ for any pair $i\ne j$.
\end{thm}

\section{Space of directions}

Recall that a metric space $\spc{X}$ is \index{geodesic space@$\ell$-geodesic space}\emph{$\ell$-geodesic}
if any two points $x,y\in\spc{X}$ such that $\dist{x}{y}{}<\ell$ can be connected by a geodesic.

\begin{thm}{Theorem}\label{thm:finite-space-of-directions}
For any point $p$ in a finite-dimensional Alexandrov space, the space of directions $\Sigma_p$ is compact and $\pi$-geodesic.
\end{thm}

By \ref{ex:GHto-tangent}, this theorem implies the following.

\begin{thm}{Corollary}\label{ex:GHto-tangent-finite-dim}
Let $p$ be a point in a finite-dimensional Alexandrov space $\spc{A}$,
and let $\lambda_n\to\infty$.
Then there is a pointed Gromov--Hausdorff convergence $(\lambda_n\cdot \spc{A},p)\z\to (\T_p,0)$.
%Moreover, for any geodesic $\gamma$ that starts at $p$, we have
%\[\iota_n\circ\gamma(t/\lambda_n)\to t\cdot \gamma^+(0),\]
%where $\iota_n$ sends a point in $\spc{A}$ to the corresponding point in $\lambda_n\cdot\spc{A}$.
\end{thm}

\parit{Proof of \ref{thm:finite-space-of-directions}.}
Choose $\eps>0$; suppose $\spc{A}$ is $m$-dimensional.
Assume $\pack_\eps \Sigma_p\ge n$;
that is, we can choose $n$ directions $\xi_1,\dots, \xi_n\in \Sigma_p$ such that $\mangle(\xi_i,\xi_j)\z>\eps$ for any $i\ne j$.
Without loss of generality, we may assume that each direction is geodesic;
that is, there is a point $x_i\in \spc{A}$ such that $\xi_i=\dir p{x_i}$.

Choose a small $r>0$ and $y_i\in [px_i]$ such that $\dist{p}{y_i}{}=r$ for each $i$.
Since $r$ is small, we can assume that $\angk p{y_i}{y_j}>\eps$ for $i\ne j$.
By \ref{cor:euclid-subcone}, we can choose $p'$ arbitrarily close to $p$ such that $\dir{p'}{y_i}\in \Lin_{p'}$ for any $i$;
in particular, all directions $\dir{p'}{y_i}$ lie in a unit sphere of dimension at most $m-1$.
Since $\dist{p'}{p}{}$ is small, $\angk {p'}{y_i}{y_j}>\eps$ if $i\ne j$.
By comparison,
\[\mangle \hinge{p'}{y_i}{y_j}>\eps.\]
Therefore, $\pack_\eps \Sigma_p\le \pack_\eps\SSS^{m-1}$.

Since $\SSS^{m-1}$ is compact, $\pack_\eps\SSS^{m-1}<\infty$ for any $\eps>0$.
By definition, the space of directions $\Sigma_p$ is complete.
Applying \ref{ex:pack-net}, we get that $\Sigma_p$ is compact.

It remains to prove the following claim.

\begin{clm}{}
If $\Sigma_p$ is compact, then it is $\pi$-geodesic.
\end{clm}

Choose two geodesic directions $\xi=\dir px$ and $\zeta=\dir py$;
let
\[\alpha\z=\tfrac12\cdot \mangle \hinge pxy=\tfrac12\cdot \dist{\xi}{\zeta}{\Sigma_p}.\]
Assume $\alpha<\tfrac\pi2$.

It is sufficient to construct an \index{almost midpoint}\emph{almost midpoint} $\dir pz$ of $\xi$ and $\zeta$ in $\Sigma_p$;
that is, we need to show that for any $\eps>0$ there is a geodesic $[pz]$ such that
\[\mangle\hinge pxz\le \alpha+\eps
\quad\text{and}\quad
\mangle\hinge pyz\le \alpha+\eps.\]
Indeed, once this is done, the compactness of $\Sigma_p$ can be used to get an actual midpoint for any two directions in $\Sigma_p$, and Menger's lemma (\ref{lem:mid>geod}) will finish the proof.

Choose a sequence of small positive numbers $r_n\to0$.
Consider sequences $x_n\z\in [px]$ and $y_n\z\in [py]$ such that $\dist{p}{x_n}{}=\dist{p}{y_n}{}=r_n$.
Let $m_n$ be a midpoint of $[x_n\,y_n]$.
Since $\alpha<\tfrac\pi2$, $m_n\ne p$.

Since $\Sigma_p$ is compact, we can pass to a subsequence of $r_n$ such that
$\dir{p}{m_n}$ converges;
denote the limit by $\mu$.
Choose a geodesic $[pz]$ that runs at a small angle to $\mu$.
It remains to show that $\dir pz$ is the needed almost midpoint.

Evidently, $\angk p{x_n}{m_n}=\angk p{y_n}{m_n}$.
By \ref{ex:alex-lemma-cat}, we have
\[\angk p{x_n}{m_n}+\angk p{y_n}{m_n}\le \angk p{x_n}{y_n}.\]

For each large $n$, choose a point $z_n\in [pz]$ such that $\dist{p}{z_n}{}=\dist{p}{m_n}{}$.
By construction, for all large $n$, we have $\mangle\hinge pz{m_n}\approx0$ with any prescribed error.
By comparison, the value $\frac{\dist{z_n}{m_n}{}}{\dist{p}{z_n}{}}$ can be assumed to be arbitrarily small for all large $n$.
Applying this observation and the definition of angle measure, we also get that the following approximations
\begin{align*}
\angk p{x_n}{y_n}&\approx \mangle\hinge p{x_n}{y_n},
\\
\angk p{x_n}{m_n}&\approx\angk p{x_n}{z_n}\approx\mangle\hinge p{x_n}{z_n},
\\
\angk p{m_n}{y_n}&\approx\angk p{z_n}{y_n}\approx\mangle\hinge p{z_n}{y_n},
\end{align*}
hold with any prescribed error for all sufficiently large $n$.
It follows that $\dir pz$ is an almost midpoint of $\dir px$ and $\dir py$, as required.
\qeds

\parit{Remark.}
In the above proof, the angles $\mangle\hinge pxz$ and $\mangle\hinge pyz$ are bounded below by
comparison, but we needed upper bounds.
These were extracted from the definition of angle measure and the compactness of the space of directions.
Halbeisen's example \cite[13.6]{alexander-kapovitch-petrunin2024} shows that it cannot be done without the compactness condition.

\begin{thm}[!]{Exercise}\label{ex:concave-differential}
Let $f$ be a semiconcave function defined in an open set $\Dom f$ in a finite-dimensional Alexandrov space.
Show that the differential $\dd_pf$ is a concave function $\T_p\to\RR$ for any $p\in \Dom f$.
\end{thm}


\section{Right-inverse theorem}

\begin{thm}{Right-inverse theorem}\label{thm:right-inverse}\index{right-inverse theorem}
Suppose $p,a_0,\dots,a_m$ are points in an Alexandrov space $\spc{A}$ such that
\[\angk p{a_i}{a_j}>\tfrac\pi2\]
for any $i\ne j$.
Then the map $f\:\spc{A}\to\RR^m$ defined by
\[f\:x\mapsto (\dist{a_1}{x}{},\dots,\dist{a_m}{x}{})\]
has a continuous right inverse $\map$ defined in a neighborhood of $f(p)$ such that $p=\map\circ f(p)$.
\end{thm}

In the proof we construct a local right inverse $\map$ of $f$ around $f(p)$.
The construction uses gradient flow for a suitably chosen family of functions.
The proof can be extracted from the following exercise;
more details are given in the hints.

\begin{thm}[!]{Exercise}\label{ex:proof-right-inverse}
Suppose $p,a_0,\dots,a_m\in\spc{A}$ and $f\:\spc{A}\to\RR^m$ are as in \ref{thm:right-inverse}.
Assume $\eps>0$ is sufficiently small.
Given $\bm{y}\z=(y_1,\z\dots,y_m)\in \RR^m$,
consider the function on $\spc{A}$ defined by
\[h_{\bm{y}}(x)\df\min\{\,0, \dist{a_1}{x}{}-y_1,\dots,\dist{a_m}{x}{}-y_m\,\}+\eps\cdot\dist{a_0}{x}{}.\]


\begin{subthm}{ex:proof-right-inverse:grad}
Show that for some fixed $r>0$ and $\lambda$, the function $h_{\bm{y}}$ is $\lambda$-concave in $\oBall(p,r)$, and the following hold at any $x\in \oBall(p,r)$:
\begin{enumerate}[(i)]
\item\label{111} $(\dd_x\distfun_{a_i})(\nabla_x h_{\bm{y}})<-\eps^2$ if $\dist{a_i}{x}{}>y_i$ and
\item\label{222} $(\dd_x\distfun_{a_i})(\nabla_x h_{\bm{y}})>\eps^2$ if
\[\dist{a_i}{x}{}-y_i=\min_{j>0}\{\dist{a_j}{x}{}\z-y_j\}<0.\]
\end{enumerate}


\end{subthm}

\begin{subthm}{ex:proof-right-inverse:alpha}
Let $\alpha_{\bm{y}}$ be an $h_{\bm{y}}$-gradient curve that starts at $p$.
Use \ref{SHORT.ex:proof-right-inverse:grad} to show that
\[
f[\alpha_{\bm{y}}(t_0)]
=
\bm{y}\]
if
$\tfrac1{\eps^2}\cdot|f(p)-\bm{y}|
\le
t_0
\le
\tfrac{r}{2}$.

\end{subthm}

\begin{subthm}{ex:proof-right-inverse:end}
Let $t_0\:\bm{y}\mapsto\tfrac{1}{\eps^2}\cdot\bigl|f(p)-\bm{y}\bigr|$.
Use \ref{lem:fg-dist-est} to show that the map
\[\map\:{\bm{y}}\mapsto \alpha_{\bm{y}}\circ t_0(\bm{y})\]
is continuous (in fact Hölder) on $\Omega=\oBall(f(p),\tfrac{\eps^2}{2}\cdot r)\subset\RR^m$
and $f\circ \Phi(\bm{y})=\bm{y}$ for any $\bm{y}\in \Omega$.
(This finishes the proof of \ref{thm:right-inverse}.)
\end{subthm}

\end{thm}

\begin{thm}{Corollary}\label{cor:tan-g-delta}
Let $\spc{A}$ be an Alexandrov space, and let $m$ be a positive integer.
If $\LinDim\spc{A}\ge m$, then
\[\dim\Lin_s\z\ge m\]
for any point $s$ in a dense G-delta subset of points in $\spc{A}$.

Moreover, if $m=\LinDim\spc{A}$, then the set $S$ of points $s\in \spc{A}$ such that $\T_s\spc{A}\iso \EE^m$ is a dense G-delta set in $\spc{A}$.
\end{thm}

It is worth mentioning that the set $S\subset\spc{A}$ is also convex; see \cite[1.10]{petrunin1998}.


\parit{Proof.}
Assume $\spc{A}$ is $\Alex0$.
By the right-inverse theorem and \ref{ex:resporka} there is a ball $B=\cBall[p,r]_{\spc{A}}$ and $\Const>0$ such that
\[\pack_\eps B\ge \frac{\Const}{\eps^m}\]
for all small $\eps>0$.

In particular, given a radius $R$ we can choose $\eps>0$ and an integer $N$ such that
\[\pack_\eps \cBall(0,R+r)_{\EE^{m-1}}< N\le \pack_\eps B.\]

Applying \ref{cor:euclid-subcone} to an $\eps$-packing $x_1,\dots,x_N$ in $B$, we get that \[\dim\Lin_s\z\ge m\]
for any point $s$ in a dense G-delta subset of $\oBall(p,R)$.
Since $R>0$ is arbitrary, the first statement follows.

To prove the last statement, apply \ref{ex:dim-max}.

The general case can be reduced to $\Alex{-1}$ using rescaling,
and it can be done essentially in the same way; one only has to change $\EE^{m-1}$ to $\HH^{m-1}$.
\qeds

\section{Distance chart}

\begin{thm}{Theorem}\label{thm:dist-chart}
Let $p,a_0,\dots,a_m$ be points in an $m$-dimensional Alexandrov space $\spc{A}$ such that
\[\angk p{a_i}{a_j}>\tfrac\pi2\]
for any $i\ne j$.
Then the map $f\:\spc{A}\to\RR^m$ defined by
\[f\:x\mapsto (\dist{a_1}{x}{},\dots,\dist{a_m}{x}{})\]
gives a bi-Lipschitz embedding of a neighborhood $\Omega$ of $p$;
the restriction $f|_\Omega$ is called a \index{distance chart}\emph{distance chart} at $p$.
\end{thm}

The following exercise provides a guide to the proof of the theorem.

\begin{thm}[!]{Exercise}\label{ex:proof-dist-chart}
Suppose $p,a_0,\dots,a_m\in\spc{A}$ and $f\:\spc{A}\to\RR^m$ are as in \ref{thm:right-inverse}.
Show that there is $\eps>0$ such that for any two points $x,y$ in a sufficiently small neighborhood of $p$,
one of the inequalities holds:
\begin{align*}
\mangle\hinge xy{a_1}&<\tfrac\pi2-\eps,\ \dots,\ \mangle\hinge xy{a_m}<\tfrac\pi2-\eps,
\\
\mangle\hinge yx{a_1}&<\tfrac\pi2-\eps,\ \dots,\ \mangle\hinge yx{a_m}<\tfrac\pi2-\eps.
\end{align*}
Use this together with the right-inverse theorem (\ref{thm:right-inverse}) to prove \ref{thm:dist-chart}.
\end{thm}


\section{Volume}

Fix a positive integer $m$.
The $m$-dimensional Hausdorff measure of a Borel set $B$ in a metric space will be called its \index{volume}\emph{$m$-volume}; it will be denoted by $\vol_m B$.
We assume that the Hausdorff measure is calibrated so that the unit cube in $\EE^m$ has unit volume.

This definition will be applied mostly to subsets of $m$-dimensional Alexandrov spaces.
In this case, we may write $\vol B$ instead of $\vol_m B$.

\begin{thm}{Bishop--Gromov inequality}\label{inq:BG}\index{Bishop--Gromov inequality}
Let $\spc{A}$ be an $\Alex0$ space.

\begin{subthm}{inq:BG:a}
For any integer $k\ge 0$, any point $p\in \spc{A}$ and any radius $r>0$, we have
\[\vol_k \oBall(p,r)_{\spc{A}}\le \vol_k \oBall(0,r)_{\T_p\spc{A}}.\]
\end{subthm}

\smallskip

\noindent Further, suppose $\LinDim \spc{A}=m<\infty$.

\smallskip

\begin{subthm}{inq:BG:b}
For any point $p\in \spc{A}$ and any radius $r>0$, we have
\[\vol_m \oBall(p,r)\le \omega_m\cdot r^m,\]
where $\omega_m$ denotes the volume of the unit ball in $\EE^m$.
\end{subthm}

\begin{subthm}{inq:BG:c}
Moreover, the function
\[r\mapsto \frac{\vol_m \oBall(p,r)}{r^m}\]
is nonincreasing.
\end{subthm}

\end{thm}

\parit{Proof.}
Given $x\in\spc{A}$, choose a geodesic path $\gamma_x$ from $p$ to $x$.
Recall that $\gamma_x^+(0)=\log_p x$ (here $\log_p\:\spc{A}\to \T_p$ is a multivalued function; so, $\gamma_x^+(0)$ might not be the only value of $\log_p x$).
By comparison, $\log_p$ is distance-noncontracting.
Note that $\log_p(\oBall(p,r)_{\spc{A}})\subset \oBall(0,r)_{\T_p}$;
this proves part \ref{SHORT.inq:BG:a}.

\begin{wrapfigure}{r}{44 mm}
\vskip-4mm
\centering
\includegraphics{mppics/pic-803}
\vskip1mm
\end{wrapfigure}

If $\T_p\iso \EE^m$, then \ref{SHORT.inq:BG:a} implies \ref{SHORT.inq:BG:b}.
In the general case, by \ref{cor:tan-g-delta}, we can find a point $p'$ arbitrarily close to $p$ such that $\T_{p'}\iso \EE^m$.
If $\eps>\dist{p}{p'}{}$, then $\oBall(p,r)\z\subset \oBall(p',r+\eps)$.
Therefore,
\[\vol \oBall(p,r)\le \omega_m\cdot (r+\eps)^m\]
for any $\eps>0$.
Hence \ref{SHORT.inq:BG:b} follows.

Now, suppose $0<r_1<r_2$.
Consider the map $w\: \spc{A}\selfmap$ defined by $w\:x\mapsto \gamma_x(\tfrac {r_1}{r_2})$; it mimics the dilation with center at $p$ and coefficient $\tfrac {r_1}{r_2}$.
By comparison,
\[\dist{w(x)}{w(y)}{}\ge \tfrac {r_1}{r_2}\cdot \dist{x}{y}{}.\]
Observe that $\oBall(p,r_1) \supset w[\oBall(p,r_2)]$.
Therefore,
\[\vol \oBall(p,r_1)\ge (\tfrac {r_1}{r_2})^m\cdot\vol \oBall(p,r_2).\]
\qedsf

The following exercise generalizes the Bishop--Gromov inequality to the $\Alex{-1}$ case;
it is sufficient for most applications.
A sharper statement for $\Alex\kappa$ spaces is given in \ref{inq:BG+}.



\begin{thm}[!]{Exercise}\label{ex:BG}
Let $\spc{A}$ be a finite-dimensional $\Alex{-1}$ space;
so $\LinDim\spc{A}=m<\infty$.
Show that
\[\vol \oBall(p,r)\le \omega_m\cdot(\sinh r)^m,\]
where $\omega_m$ denotes the volume of the unit ball in $\EE^m$.
Moreover, the function
\[r\mapsto \frac{\vol \oBall(p,r)}{(\sinh r)^m}\]
is nonincreasing.
\end{thm}

\begin{thm}[!]{Exercise}\label{ex:diam-compact:proper}
Show that any finite-dimensional Alexandrov space is proper.

\end{thm}

\begin{thm}{Exercise}\label{ex:vol>0}
Show that any nonempty open set in an Alexandrov space of finite dimension $m$ has positive $m$-volume.
\end{thm}


\section{Packings}\label{sec:packing-numbers}


\begin{thm}{Exercise}\label{ex:tangent=Em}
Let $\spc{A}$ be an $m$-dimensional $\Alex\kappa$ space and $m\z<\infty$.

\begin{subthm}{ex:tangent=Em:2balls}
Show that for any two balls $B_1$ and $B_2$ in $\spc{A}$ with positive radii,
there is $\delta>0$ such that the inequality
$\pack_{\delta\cdot\eps} B_1\ge \pack_\eps B_2$
holds for all sufficiently small $\eps>0$.
\end{subthm}

\begin{subthm}{ex:tangent=Em:ball}
Show that there are constants $c_1=c_1(m)$ and $c_2=c_2(m,r,\kappa)$ such that
\[c_1\cdot \vol B\le\eps^m\cdot\pack_\eps B\le c_2\cdot \vol B\]
for any ball $B=\cBall[p,r]_{\spc{A}}$ with $r>0$ and all small $\eps>0$.
\end{subthm}

\end{thm}


The exercise implies that for Alexandrov spaces, \index{Minkowski dimension}\emph{Minkowski dimension} (also known as \index{box-counting dimension}\emph{box-counting dimension}) coincides with linear dimension.
Namely, let $\spc{A}$ be an Alexandrov space.
Given a point $p\in \spc{A}$ and $r>0$, let
\[m(p,r)\df \inf\set{\alpha\in\RR}{\eps^\alpha\cdot\pack_\eps \cBall[p,r]\to 0\quad\text{as}\quad \eps\to 0}.\]
Then $m(p,r)=\LinDim \spc{A}$; in particular, $m(p,r)$ is an integer, and it does not depend on the choice of $p$ and $r$.


\section{Other dimensions}\label{sec:all-dim}

Now we want to show that \textit{all reasonable definitions of dimension give the same result for Alexandrov spaces}.
More precisely, we have the following theorem; compare with \cite[15.16]{alexander-kapovitch-petrunin2024}.
We refer to \cite{hurewicz-wallman} for definitions of \index{Lebesgue covering dimension}\emph{Lebesgue covering dimension} and \index{Hausdorff!dimension}\emph{Hausdorff dimension}, which will be denoted by \index{TopDim@$\TopDim$}$\TopDim$ and
\index{HausDim@$\HausDim$}$\HausDim$, respectively.

\begin{thm}{Theorem}\label{thm:dim=dim}
For any Alexandrov space $\spc{A}$, we have
\[\LinDim \spc{A}=\TopDim \spc{A}=\HausDim \spc{A}.\]
\end{thm}

\parit{Proof.}
Suppose $\LinDim \spc{A}=\infty$.
By the right-inverse theorem (\ref{thm:right-inverse}), $\spc{A}$ contains a compact subset $K$ with arbitrarily large $\TopDim K$.
In particular,
\[\TopDim\spc{A}=\infty.\]
By Szpilrajn's theorem,
$\HausDim K\ge \TopDim K$.
Thus we also have
\[\HausDim\spc{A}=\infty.\]

Now suppose $\LinDim \spc{A}=m<\infty$.
By the Bishop--Gromov inequality (\ref{inq:BG} and \ref{ex:BG}),
\[\HausDim\spc{A}\le m.\]

Since $\spc{A}$ is proper (\ref{ex:diam-compact:proper}),
Szpilrajn's theorem implies that
\[\TopDim\spc{A}\le \HausDim\spc{A}\le m.\]
Finally, by the right-inverse theorem (\ref{thm:right-inverse}), $m\le\TopDim\spc{A}$.
\qeds

\begin{thm}{Exercise}\label{ex:dim=dim}
Let $\Omega$ be an open subset of an Alexandrov space $\spc{A}$.
Show that
\[\LinDim \spc{A}=\LinDim \Omega=\TopDim \Omega=\HausDim \Omega.\]
\end{thm}

\begin{thm}[!]{Exercise}\label{ex:finite-tan}
Let $p$ be a point in a finite-dimensional Alexandrov space $\spc{A}$.
Prove the following:
\begin{subthm}{ex:finite-tan:tan}
The tangent space $\T_p$ is a proper $\Alex0$ space.
\end{subthm}

\begin{subthm}{ex:finite-space-of-directions-dim}
$\LinDim\Sigma_p+1=\LinDim\T_p=\LinDim\spc{A}$.
\end{subthm}

\begin{subthm}{ex:finite-tan:sigma}
If $\LinDim \spc{A}>1$, then $\Sigma_p$ is geodesic.
\end{subthm}

\end{thm}

Part \ref{SHORT.ex:finite-space-of-directions-dim} of the exercise opens a way to use induction on dimension, which can be used for finite-dimensional Alexandrov spaces;
in Lecture~\ref{chap:bry}, we will use it to define boundary.
The condition $\LinDim \spc{A}>1$ in \ref{SHORT.ex:finite-tan:sigma} is necessary;
indeed $\Sigma_0\RR$ is a two-point space which is not geodesic.

\section{Remarks}

Halbeisen's example \cite{alexander-kapovitch-petrunin2024,halbeisen} shows that compactness of the space of directions is essential in the proof that the space of directions is $\pi$-geodesic (\ref{thm:finite-space-of-directions}).

The proof of the following version of the Bishop--Gromov inequality additionally requires the so-called \textit{coarea formula} for Alexandrov spaces.
The weaker inequality from \ref{ex:BG} is sufficient for the sequel.

\begin{thm}{Optimal Bishop--Gromov inequality}\label{inq:BG+}\index{Bishop--Gromov inequality}
Given a point $p$ in an $m$-dimensional $\Alex\kappa$ space,
consider the function
\[v\:r\mapsto\vol_m\oBall(p,r);\]
denote by $\tilde v(r)$ the volume of an $r$-ball in $\MM^m(\kappa)$.
Then
\[v(r)\le \tilde v(r)\]
for any $r>0$.
Moreover, the function
\[r\mapsto \frac{v(r)}{\tilde v(r)}\] is nonincreasing.
If $\kappa>0$, then one has to assume that $r<\tfrac\pi{\sqrt\kappa}$.
\end{thm}

This inequality was originally proved for Riemannian manifolds with Ricci curvature bounded below.
The first part is also known as \index{Bishop's inequality}\emph{Bishop's inequality}; it is due to Richard Bishop \cite{bishop1964}, \cite[Corollary 4, p. 256]{bishop-crittenden}.
The second part is due to Michael Gromov \cite{gromov1981}.

Theorem~\ref{thm:dim=dim} was essentially proved by Conrad Plaut \cite{plaut:dimension}.
At that time, it was not known whether
\[\LinDim\spc{A}=\infty\quad \Rightarrow\quad \TopDim\spc{A}=\infty\]
for any Alexandrov space $\spc{A}$.
The latter implication was proved by Grigory Perelman and the third author \cite{perelman-petrunin:qg}; our proof is slightly different.
